% ============================================================================
%  Subdocumento 2. Contenido. Se usa igual en el documento único y en el
%  PDF propio del subdocumento. Los formularios que le asigna el T-21 se
%  agregan solos al final (datos en formularios/ de esta carpeta).
% ============================================================================

\begin{resumenApertura}
\marcador{Qué resuelve este subdocumento, en cuatro o cinco líneas}
\begin{recibe}
  \item \marcador{Qué recibe el mandante}
\end{recibe}
\end{resumenApertura}

\section{Resumen ejecutivo}\label{sec:2-resumen-ejecutivo}

\marcador{La propuesta de valor, el enfoque de ingeniería y los resultados comprometidos}

\section{Dimensionamiento del problema}\label{sec:2-dimensionamiento-del-problema}

\marcador{Magnitud del problema con datos cuantitativos del Caso, jerarquía de patologías y orden de resolución}

\section{Contexto de la industria, regulación y estacionalidad}\label{sec:2-contexto-de-la-industria-regulac}

\marcador{Particularidades operacionales, regulatorias y estacionales, y el control que exige cada norma}

\section{Actores y grupos de interés}\label{sec:2-actores-y-grupos-de-interes}

\marcador{Actores afectados, influencia e interés, y la posición de audIT frente a cada tensión}

\section{Supuestos declarados}\label{sec:2-supuestos-declarados}

\marcador{Supuestos propios de la propuesta con fundamento e impacto si resultan falsos}
